\section{Eleven positive variables: transport only the two small gaps}
\label{sec:positive211}
The new polynomial has eleven positive variables, degree $211$, and $113$
live binary arithmetic operations.  It improves the incoming centered
$257/114$ construction in both degree and evaluation cost.  The earlier
$277/111$ point still uses two fewer operations.  Centering itself is not
new: our contribution is a second center, a cheaper gap bound, and a
cofinality argument that permits a useful change of the auxiliary witness.

\subsection{The literal polynomial}
Use the complete compact44 coefficient tuple
$\beta,K_0,M_C,M_F$, with $M_F$ already shifted, as bound in
\texttt{positive/native/ACTIVE\_BINDING.json}.  The unknowns are
\[
 j,F,\alpha,i,T,X,Y,n,\eta,\zeta,\gamma\in\mathbb Z_{>0}.
\]
Set
\begin{align*}
 J&=3j+1,&q&=\beta J+1,&C&=q-F-\alpha,\\
 E&=XY,&k&=\eta+\zeta,&c&=kY+\eta,\\
 a&=E+Y,&H&=4a+3,&\Delta&=a^2+H=(a+2)^2-1,\\
 \ell&=EkY,&I&=ic^2,&D_f&=1+\Delta I^2,\\
 Q&=\Delta(D_f-1),&\delta&=Q-1,&D_\tau&=\ell(\ell+k)+1.
\end{align*}
The coefficient and transport registers are
\begin{align*}
 R&=((q-1)(q-F)+\alpha)(q^2-1)+(M_C+qM_F)J,\\
 M&=q^3(q-1),\qquad B_t=(qK_0+X)C-qF.
\end{align*}
Here are the new centers and the small transport gap:
\begin{align*}
 T_o&=T+q+\alpha,& b&=cT_o-R,&V&=D_f b-c,\\
 D&=(a+2)I,&L&=kT_o+F,&W&=1+(q-1)L,\\
 \overline Z&=M(\ell+k+n+L+V+D)+(q-1)Y-q^3WB_t.
\end{align*}
Keep the adaptive main carrier
\begin{align*}
 r_0&=k-R-1,&S&=4qY,&U&=qS,\\
 D_*&=U(X+ac)+H(U\gamma+SX-r_0),\qquad
 A_1=D_*^2-U^2(1+\Delta c^2).
\end{align*}
Define the three-root norm without any division:
\begin{align*}
 A&=M^2D_\tau,\quad B=M^2D_f,\quad C_N=M^2(QV^2-1),\\
 T_N&=\delta^2AB,\quad R_N=\delta C_N,\\
 J_N&=\delta(\overline Z^2-A-B)-C_N,\\
 N_c&=J_N^2-4[T_N+R_N(A+B)],\\
 N_8&=N_c^2-64T_NR_N\overline Z^2,\\
 \boxed{P_{11}}&\boxed{=N_8+A_1+N_8A_1=0.}
\end{align*}
All registers are abbreviations, not additional quantified coordinates.
\par The exact circuit is \path{positive/p11_degree211.json}.

\subsection{Recover divisibility before using positivity}
Write, only in the proof,
\[
 D_y=\frac{QV^2-1}{Q-1},\qquad
 \tau=\sqrt{D_\tau},\quad f=\sqrt{D_f},\quad y=\sqrt{D_y}.
\]
The inherited unconditional separation lemma applies for every integral
$b$, without assuming $C>0$ or $V>0$.  Indeed $I\ge c^2$, $c\ge3$ imply
$1<D_\tau<D_f<Q$ and $|V|>1$.  If $b=0$, then $V=-c$ and
$f>\tau+y$; otherwise $y>f+\tau$.  The radicands are distinct, since
$D_y=D\in\{D_\tau,D_f\}$ would imply $Q(D-V^2)=D-1$, with
$0<D-1<Q$.  Thus no signed sum vanishes.

The norm equation gives
$\overline Z/M=\tau_*+f_*+y_*$ with the appropriate signed roots.
These roots are rational.  An explicit extraction identity is as follows:
for squares $a_1,a_2,a_3$ and nonzero sum $s=x_1+x_2+x_3$, put
\[
 u=\frac{s^2-a_1-a_2-a_3}{2},\qquad
 v=\frac{u^2-a_1a_2-a_1a_3-a_2a_3}{2s}.
\]
Then $x_i=a_i(a_i+u)/(a_is+v)$; the denominator equals
$x_i(x_i+x_j)(x_i+x_k)$ and cannot vanish by distinctness.
The first two roots are integral because their squares are integers.
The adjacent-Pell rational-integrality lemma makes $y_*$ integral from
$QV^2-(Q-1)y_*^2=1$.  For completeness, write $y_*=a/m$ reduced;
$m^2\mid Q-1$, and the norm-preserving descent
\[
 V'=(2Q-1)V-2(Q-1)y_*,\qquad
 y'=(2Q-1)y_*-2QV
\]
has integral $V'$, unchanged reduced denominator $m$, and
$0<|V'|<|V|$ when $|V|>1$.  The terminal case $|V|=1$ forces $m=1$.

Therefore $M\mid\overline Z$, and reduction modulo $q^3$ gives
$q^3\mid Y$.  This is the key dependency order: no transport-gap sign
or positivity of $C$ has been used yet.

\subsection{A degree-two bound for the auxiliary gap}
The authentic masks give $0<(M_C+qM_F)J<2q^2$ and $q\ge16$.  Directly,
\[
 -Fq^3<R<q^3(q+\alpha).
\]
For the upper inequality, subtract the bound
$q(q-1)(q^2-1)+\alpha(q^2-1)+2q^2$ from $q^3(q+\alpha)$;
the difference is $q^3-q^2-q+\alpha(q^3-q^2+1)>0$.
Since $c>Y\ge q^3$, we get
\[
 c(q+\alpha)>R,\qquad b>cT\ge c,\qquad V>1.
\]
Also $\delta>(\Delta-1)D_f$, $c<2kY$, and
$\Delta-1=(a+2)^2-2>4Y^2$.  Consequently
\begin{align*}
 0<y-V
 &=\frac{V^2-1}{\delta(y+V)}
 <\frac{b}{\Delta-1}\\
 &<\frac{cT_o+Fq^3}{\Delta-1}<kT_o+F=L.
\end{align*}
The second center satisfies
\[
 D^2-D_f=I^2-1>0,\qquad 0<D-f=\frac{I^2-1}{D+f}<I.
\]
Thus the huge root $y$ is transported only through $y-V$, and the strong
root $f$ through the favorable gap $D-f$.

\subsection{Soundness and cofinal completeness}
Recover positive witnesses
\[
 u=y-V,\qquad
 r=n+L+V+D-f_*-y_*,\qquad f_{\rm old}=|f_*|,\quad i_{\rm old}=i.
\]
Indeed $r\ge n+L-(y-V)+D-f>0$.  The weighted first carrier with
this $r$ and weight $L$ is $M\tau_*$.  The strong and auxiliary residuals,
and $A_1$, vanish.  The inherited weighted-source theorem therefore yields
a finite ZFC refutation.  The signed roots, not their absolute values,
must be used inside $r$.

Conversely, canonical source witnesses give the outer coordinates and a
normalized strong pair $f,I$, and the inherited auxiliary lemma supplies
one positive pair $V_0,y_0$ with
\[
 QV_0^2-(Q-1)y_0^2=1,\qquad V_0=f^2(cT_0-R)-c.
\]
The matrix
\[
 B_Q=\begin{pmatrix}2Q-1&2(Q-1)\\2Q&2Q-1\end{pmatrix}
\]
has determinant one, positive entries, and preserves the norm.
Some positive power is the identity modulo $h=cf^2$, by finiteness of
$\operatorname{SL}_2(\mathbb Z/h\mathbb Z)$.  Its iterates give unbounded
positive pairs with $V_j\equiv V_0\pmod h$, so
\[
 T_j=(V_j+c+Rf^2)/(cf^2)
\]
is integral and unbounded.  Choose $T_o=T_j>q+\alpha$ with
$L=kT_o+F>I$.  Then $T=T_o-q-\alpha>0$.
The canonical weighted carrier supplies
$r>qW-2k>2L>L+I$; hence
\[
 n=r-L-V-D+f+y=r-L-(D-f)+(y-V)>0.
\]
It has $\overline Z=M(\tau+f+y)$, and the unchanged adaptive carrier
supplies positive $\gamma$.  This proves completeness without assuming
an unjustified size condition on the original auxiliary witness.
Finally, $N_8$ and $A_1$ are squares modulo four, so their unit-product
finalizer vanishes exactly when both vanish.  Negating solvability gives
$\operatorname{Con}(\mathrm{ZFC})$ in ZFC, with the same explicitly
inherited compiler and calibration boundary.

\subsection{Exact degree, not a generic-gate bound}
Put $z=\overline Z-M(V+D)$, $e=I^2-1$, and
$H_0=z^2+2MDz+M^2e-A$.  Then
\[
 J_N=\delta[H_0+2MV(z+MD)]-M^2(V^2-1).
\]
The identity $B=M^2(D^2-e)$ cancels the apparent degree-102 part of
$N_c$.  Its degree-99 head is $8Q_0^2M_0^3V_0^2D_0z_0$, and the
norm has degree-198 head
\[
 64Q_0^4M_0^6V_0^4D_0^2(z_0^2-M_0^2D_{\tau,0}).
\]
With $q_0=3\beta j$, $k_0=\eta+\zeta$,
$L_0=k_0(T+q_0+\alpha)$ and
$B_{t,0}=(K_0q_0+X)(q_0-F-\alpha)-q_0F$, we have
\[
 z_0=M_0(\ell_0-L_0B_{t,0}),\qquad
 z_0^2-M_0^2D_{\tau,0}
 =M_0^2L_0B_{t,0}(L_0B_{t,0}-2\ell_0)\ne0.
\]
Multiplying by the degree-13 adaptive-carrier head shows that the
coefficient of
\[
 j^{48}i^{18}\eta^{39}X^{30}Y^{67}F^4T^2\alpha^2\gamma
\]
is $8192(3\beta)^{48}>0$.  This certifies exact degree 211 on the
actual coefficient slice, independently of $K_0,M_C,M_F$.

\subsection{A materialized incoming degree70 companion}
\sloppy
The manuscript-only incoming \texttt{report.zip} already proposes a
12-positive-variable equation of degree 70 with 100 operations.  Its advertised
JSON was absent.  We supply it as
\texttt{positive/p12\_degree70.json}, with independent checks matching
all metrics, but do not claim its theorem as new here.  It replaces the
normalized strong norm by the ordinary strong norm after an early index
bound licenses ordinary rank recovery.  Its exact residual degrees are
$15,13,42$.  The full proof is retained verbatim in
\path{positive/interfaces/incoming_report.tex}; the crucial endpoint
is the ordinary source, not a possibly nonintegral division of the strong
multiplier by $\Delta$.
